\section{Propeller Spinner}

File \texttt{propeller\_spinner.py} không phải là PID controller, nhưng giúp
cánh quạt trong Isaac Sim quay theo RPM của motor model. Module này dùng
\texttt{dynamic\_control} để tìm các DOF:

\begin{equation}
    propeller\_1\_joint,\quad
    propeller\_2\_joint,\quad
    propeller\_3\_joint,\quad
    propeller\_4\_joint
\end{equation}

Nếu tìm được DOF, code cấu hình velocity drive và đặt target velocity cho từng
cánh quạt.

\subsection{Đổi RPM sang vận tốc joint}

Với motor thứ $i$, RPM từ motor model được đổi sang rad/s:

\begin{equation}
    \omega_i
    =
    RPM_i
    \frac{2\pi}{60}
\end{equation}

Lệnh vận tốc joint:

\begin{equation}
    \omega_{cmd,i}
    =
    d_i
    s
    \omega_i
\end{equation}

Trong đó:

\begin{itemize}
    \item $d_i$ là chiều quay motor, mặc định $[1,-1,1,-1]$.
    \item $s$ là hệ số hiển thị, mặc định $0.7$.
\end{itemize}

Việc nhân với $s$ giúp tốc độ quay hiển thị không quá cực đoan so với mô phỏng.
Phần này chủ yếu phục vụ trực quan, không phải nguồn lực chính tác động lên
drone. Lực chính vẫn đến từ \texttt{MotorModel} và
\texttt{apply\_body\_force}.

\section{Tracking Logger}

File \texttt{tracking\_logger.py} ghi log trạng thái drone và target ra CSV.
Logger được bật bằng:

\begin{equation}
    start\_tracking\_log(csv\_path,\ sample\_period)
\end{equation}

Mặc định:

\begin{equation}
    sample\_period = 0.02\ \mathrm{s}
\end{equation}

Tức là tần số ghi log tối đa khoảng:

\begin{equation}
    50\ \mathrm{Hz}
\end{equation}

\subsection{Các đại lượng được ghi}

CSV log gồm các nhóm dữ liệu:

\begin{itemize}
    \item Thời gian mô phỏng: \texttt{time}.
    \item Vị trí và target: $x,x_d,y,y_d,z,z_d$.
    \item Roll, pitch, yaw theo độ.
    \item Roll, pitch, yaw target theo độ.
    \item Roll, pitch, yaw theo radian.
    \item Roll, pitch, yaw target theo radian.
\end{itemize}

Các dữ liệu này đủ để vẽ đồ thị tracking, kiểm tra overshoot, kiểm tra sai số
xác lập và so sánh target với trạng thái thật.

\section{Compatibility Mixin}

File \texttt{attitude\_hold\_compat.py} chứa
\texttt{AttitudeHoldCompatibilityMixin}. Module này không tạo luật điều khiển
mới. Vai trò của nó là giữ lại API cũ sau khi code được tách thành nhiều
controller nhỏ.

Ví dụ, khi code bên ngoài truy cập:

\begin{equation}
    kp\_x
\end{equation}

thì property trong mixin thực chất đọc:

\begin{equation}
    position\_controller.kp\_x
\end{equation}

Tương tự:

\begin{equation}
    z\_target
    \leftrightarrow
    altitude\_controller.z\_target
\end{equation}

\begin{equation}
    yaw\_target
    \leftrightarrow
    yaw\_controller.yaw\_target
\end{equation}

\begin{equation}
    roll\_target,\ pitch\_target
    \leftrightarrow
    attitude\_controller
\end{equation}

Cách làm này giúp những script cũ vẫn chạy được, đồng thời code mới vẫn có cấu
trúc rõ ràng hơn.

\section{Tóm tắt toàn bộ hệ thống}

Toàn bộ controller có thể tóm tắt bằng chuỗi biến đổi:

\begin{equation}
    (x_d,y_d,z_d,\psi_d)
    -
    (x,y,z,\psi)
    \longrightarrow
    PID
    \longrightarrow
    PWM_i
    \longrightarrow
    T_i,Q_i
    \longrightarrow
    \text{Isaac Sim physics}
\end{equation}

Trong đó:

\begin{itemize}
    \item Position PID biến sai số $x,y$ thành góc nghiêng mong muốn.
    \item Altitude PID biến sai số $z$ thành PWM nền.
    \item Attitude PID biến sai số roll/pitch thành PWM correction.
    \item Yaw PID biến sai số yaw thành PWM correction.
    \item Mixer tạo PWM cho từng motor theo cấu hình quad X.
    \item Motor model biến PWM thành lực đẩy và torque phản lực.
    \item Isaac Sim nhận lực vật lý và cập nhật trạng thái drone.
\end{itemize}

Điểm quan trọng nhất là mỗi PID chỉ tạo output mà tầng sau cần. Nhờ vậy hệ điều
khiển dễ đọc, dễ chỉnh tham số và dễ debug: khi drone lỗi độ cao thì kiểm tra
altitude PID; khi nghiêng sai thì kiểm tra attitude PID/mixer; khi trôi vị trí
ngang thì kiểm tra position PID và quy ước dấu roll/pitch.
